\chapter{Guida alla lettura}

\section*{Organizzazione}

Le note seguono il programma del corso, ma gli argomenti sono disposti in modo che ogni concetto venga
\textbf{definito una sola volta}, nel punto in cui serve per la prima volta, e poi soltanto richiamato.

\begin{itemize}
  \item \textbf{Parte I --- Fondamenti.} Strumenti matematici, formulazione del problema dell'apprendimento
        (dati, task, modello, algoritmo, loss) e ruolo del bias induttivo.
  \item \textbf{Parte II --- Primi modelli.} Modelli lineari (con la discesa del gradiente), overfitting e
        regolarizzazione, K-nearest neighbors: i due estremi \emph{rigido} e \emph{flessibile}.
  \item \textbf{Parte III --- Generalizzazione.} La teoria (VC-dimension, bound, SRM), la decomposizione
        bias-varianza con gli ensemble, e la stima empirica dell'errore (validazione).
  \item \textbf{Parte IV --- Reti neurali.} Perceptron, MLP, backpropagation e addestramento in pratica.
  \item \textbf{Parte V --- SVM.} Margine massimo, soft margin, kernel, regressione, metodi kernel.
  \item \textbf{Parte VI --- Modelli avanzati.} CNN, deep learning, reti randomizzate, apprendimento non
        supervisionato, reti ricorrenti, grafi.
\end{itemize}

Ogni capitolo si apre con un breve riquadro che ne descrive il contenuto e si chiude con i \emph{punti chiave}
e alcune \emph{domande di verifica} (tratte dalle domande d'esame indicate nel corso).

\section*{Convenzioni grafiche}

\begin{definizione}[nome]
Definizioni formali dei concetti.
\end{definizione}
\begin{teorema}[nome]
Enunciati di teoremi e risultati; la dimostrazione, quando prevista, segue subito dopo il riquadro.
\end{teorema}
\begin{algoritmo}[nome]
Procedure passo per passo.
\end{algoritmo}
\begin{idea}
Il concetto che bisogna aver capito di una sezione.
\end{idea}
\begin{attenzione}
Errori frequenti e affermazioni da non fraintendere. È l'unico riquadro colorato.
\end{attenzione}
\begin{esempio}
Esempi ed esercizi svolti.
\end{esempio}

Nel testo il \textbf{grassetto} segnala un termine definito per la prima volta, il \ingl{corsivo} i termini
inglesi e le enfasi, il \uline{sottolineato} le poche affermazioni da ricordare alla lettera. Le formule
riquadrate sono quelle da saper scrivere a memoria.

\section*{Notazione}

\begin{center}
\small
\begin{tabularx}{0.92\linewidth}{>{\raggedright\arraybackslash}p{3.2cm}X}
\toprule
Simbolo & Significato \\
\midrule
$x$, $\x$, $X$ & scalare, vettore (colonna, in grassetto), matrice \\
$\x\T\vv{y}$ & prodotto scalare; $\T$ indica la trasposta \\
$l$ & numero di esempi del dataset, in tutti i capitoli \\
$n$ & dimensione dell'input, cioè numero di feature \\
$\x_p$, $x_{p,i}$ & $p$-esimo esempio (riga del dataset), sua componente $i$ \\
$d_p$ & target dell'esempio $p$ (da \ingl{desired}, la risposta desiderata) \\
$o$ & uscita di un'unità o della rete \\
$h$, $H$ & ipotesi, spazio delle ipotesi \\
$\h$ & stato nascosto (strato nascosto, stato di una rete ricorrente o di un nodo di un grafo) \\
$f$ & funzione target (sconosciuta); nelle reti: funzione di attivazione \\
$\w$, $w_0$ & pesi (parametri liberi), bias (indicato anche con $b$) \\
$L$, $E$ & loss sul singolo esempio; errore (somma o media delle loss, come indicato caso per caso; con il fattore
$\frac12$ nella backpropagation) \\
$R$, $\Remp$ & rischio (errore vero), rischio empirico \\
$K$, $k$ & numero di gruppi (fold, cluster, classi, modelli di un comitato); numero di vicini nel K-NN \\
$\eta$, $\lambda$, $\alpha$ & learning rate, coefficiente di regolarizzazione, momentum (nelle SVM $\alpha_i$ sono i
moltiplicatori di Lagrange) \\
TR, VL, TS & training set, validation set, test set \\
\bottomrule
\end{tabularx}
\end{center}
